\section{Design decisions: storage and concurrency}
\label{sec:decisions-storage}

\subsection{Storage selection}

Stage~1 compared several structures for each store and chose, for 200 books, a datalake with one folder per book and a
monolithic JSON inverted index. Stage~2 must index tens of thousands of books with several processes sharing the data,
so the group moved to structures that Stage~1 had already implemented and tested:

\begin{table}[H]
\centering\small
\begin{tabularx}{\linewidth}{@{}L{2.2cm}L{3.3cm}X@{}}
\toprule
\textbf{Store} & \textbf{Choice} & \textbf{Why} \\
\midrule
Datalake & \emph{range}: one folder per 1,000 ids, such as \code{01000-01999/} &
  With one folder per book the root holds one entry per book, hundreds of thousands. With \emph{range} each folder has
  at most about 3,000 files (header, body and \code{\_READY} of each book), and the folder of a book is computed from
  its id, never searched. Files are the cheapest way to keep raw text, and the only store on disk besides MongoDB. \\
Inverted index & MongoDB: one document \code{\{term, postings\}} per term, unique index on \code{term} &
  The monolithic JSON is held whole in memory and rewritten whole at each flush; Stage~1 itself noted that it had to
  change beyond about 2,000 books. In MongoDB a flush sends only the terms that changed, a query reads only the
  documents of its terms, and many processes can read and write it at once. \\
Metadata & MongoDB: collection \code{books}, one document per book &
  By volume SQLite was enough (about 30~MB for 100,000 books). The reason is Stage~3: a SQLite file cannot be shared
  between machines, and MongoDB was already deployed. Search reads the metadata by \code{\_id} and filters in memory. \\
Publication & MongoDB: four collections &
  \code{index\_meta}, \code{index\_books}, \code{index\_publications} and \code{index\_log}. The version of the index must change for Search in one atomic step, and a conditional update of one document
  (\code{CURRENT}) gives it without any lock (Section~\ref{sec:indexing}). \\
\bottomrule
\end{tabularx}
\caption{Storage of each kind of data. The datalake is the only persistent truth on disk; everything else is in one
MongoDB database.}
\label{tab:storage}
\end{table}

\textbf{What the measures say about this choice} (Section~\ref{sec:benchmarks-indexing}). With 30,000 books the index has
651,565 terms and 94.7 million ids, and takes 1.11~GB on disk. Its longest posting list, the word \emph{said}, holds
29,898 ids in 318~KB: far from the 16~MB limit of a MongoDB document, and reading it takes 2.8~ms. One document per
term is therefore enough for queries at this scale. The cost is on the write side. \code{\$addToSet} compares each new
id with the whole list to keep it free of duplicates, so adding 100 ids costs 6~ms on a list of 1,000 ids and 233~ms on
one of 100,000. Since the frequent terms hold almost one id per book, every batch pays for the size of the index, and
the indexing rate falls from about 17 to 1 book per second between an empty index and 30,000 books.

\textbf{What would fix it, and why it was not done.} \code{\$push} is 14 to 20 times faster on long lists, but it accepts
duplicates, and idempotency is what makes every retry of the pipeline safe (Section~\ref{sec:indexing}). Avoiding
duplicates without \code{\$addToSet} would need another mechanism, such as writing a batch only once per book with a
check in \code{index\_log}. The other option is to split each posting list into buckets of bounded size, which keeps
\code{\$addToSet} cheap but changes the contract of the index and the reads of Search. Both are proposed for Stage~3, not
measured.

\subsection{Concurrency model of Indexing}

One Indexing process runs these threads:

\begin{table}[H]
\centering\small
\begin{tabularx}{\linewidth}{@{}L{3.6cm}L{2.2cm}X@{}}
\toprule
\textbf{Threads} & \textbf{How many} & \textbf{What they do} \\
\midrule
Jetty pool & up to \code{INDEXING\_HTTP\_THREADS} (250) & One per HTTP request \\
Parsers & \code{INDEXING\_THREADS} (4) & Read, parse and tokenize the books of a batch, in parallel \\
Queue writer & 1 & Holds \code{index-writer}, claims, writes, completes \\
Lease renewer & 1 & Renews the leases of the batch being written \\
Notifier & 1 & Announces each publication and marks its books \code{PUBLISHED} \\
Rebuild & 1, and 1 for its marker & Only while a rebuild runs \\
\bottomrule
\end{tabularx}
\caption{Threads of one Indexing process. The Hazelcast client and the MongoDB driver add their own.}
\end{table}

\textbf{Many parsers, one writer.} Preparing a book is pure CPU and independent of the others, so it runs in parallel.
Writing is not: the pending changes of the inverted index are one map per process, and the order of the writes
(metadata, posting lists, \code{index\_log}) is what makes a book complete. So writes go one at a time, whether they come
from the queue writer or from \code{POST /index/update}. Between replicas, the \code{index-writer} lock leaves a single
writer in the whole system; more replicas are standbys for failover, not more throughput. A second writer would update
the same documents as the first, because the frequent terms are in almost every book: the two would contend on the same
posting lists in MongoDB and add no capacity where it is missing.

\textbf{Protection of the shared state.}

\begin{table}[H]
\centering\small
\begin{tabularx}{\linewidth}{@{}L{3.6cm}L{3.8cm}X@{}}
\toprule
\textbf{Shared state} & \textbf{Shared by} & \textbf{Protection} \\
\midrule
Pending posting lists & Threads of one process & A lock around add and flush \\
Writing the index & Replicas of Indexing & Hazelcast lock \code{index-writer}, held by one thread \\
Leases of the claimed books & Writer and Control & Conditional replace of the task in \code{book-state} \\
\code{CURRENT} & Replicas that publish & Conditional update: only from V to V+1 \\
Version of a published book & Publications & \code{\$setOnInsert}: it never changes \\
Posting lists and metadata & Replicas, retries & Idempotent writes: \code{\$addToSet}, upserts by \code{\_id} \\
Rebuild marker & Replicas & \code{putIfAbsent} with a TTL and a token; renewed and removed only by its owner \\
\bottomrule
\end{tabularx}
\caption{How Indexing protects each piece of shared state.}
\end{table}

\subsection{Where threads and replicas help, and where they contend}

\begin{itemize}
  \item \textbf{More parser threads help the CPU part only.} Alone, preparing a batch goes from 233 to 487 books per
        second between 1 and 8 threads, and stops improving at 4. Inside the deployment the whole path into MongoDB
        stays at 15 books per second whatever the threads: CPU is not what limits indexing.
  \item \textbf{More Indexing replicas do not add throughput} by design: they are standbys. What they add is
        failover, measured at 2.8 to 4.0~s.
  \item \textbf{MongoDB is the shared resource.} Indexing writes the index that Search reads. With the pipeline
        running, Indexing processed 2.42 books per second alone and 0.91 while Search was under load, and with 128
        clients the p95 of Search rose from 109 to 160~ms. On one host they also share its CPUs.
  \item \textbf{What would help}, for Stage~3: a cheaper write than \code{\$addToSet} on long lists, a larger MongoDB
        cache (the index exceeds its 0.25~GB from about 10,000 books), and MongoDB on a machine of its own, away from
        the other services.
\end{itemize}
